%%%%%%%%%%%%%%%
% 导言区配置 / preamble configuration
% 由 cv-main.tex 载入：%%%%%%%%%%%%%%%
% 导言区配置 / preamble configuration
% 由 cv-main.tex 载入：%%%%%%%%%%%%%%%
% 导言区配置 / preamble configuration
% 由 cv-main.tex 载入：%%%%%%%%%%%%%%%
% 导言区配置 / preamble configuration
% 由 cv-main.tex 载入：\input{preamble}
% 使用说明见 README.md
%
% 载入时机：必须在 \usepackage{settings} 之后（依赖其加载的 biblatex）
% 例外：\PassOptionsToPackage{...}{biblatex} 必须留在 cv-main.tex 里
%       （它得在 settings.sty 载入 biblatex 之前生效）
%

% 去掉 URL 前的 "[Online] Available:" 前缀
\DefineBibliographyStrings{english}{url={\textsc{url}}}

%% 自定义文献排序：年份降序 → 本人作者位次(sortkey)升序 → 作者 → 标题
%% sortkey 由 sections/04_publications.bib 提供，1 = 第一作者
%% 详见 README「参考文献：排序」
\DeclareSortingTemplate{ryc}{
  \sort{\field{presort}}
  \sort[direction=descending]{\field{sortyear}\field{year}\literal{9999}}
  \sort{\field{sortkey}\literal{99}}
  \sort{\field{sortname}\field{author}\field{editor}\field{translator}\field{sorttitle}\field{title}}
  \sort{\field{sorttitle}\field{title}}
}
\ExecuteBibliographyOptions{sorting=ryc}

%% 作者列表截断到本人为止，如 "C. 7au et al."、"Smith, C. 7au and others"
%% sortkey 是 skipout 字段(不写入 .bbl)，故先复制到通用字段 usera 供排版使用
%% 详见 README「作者列表截断」
\DeclareSourcemap{
  \maps[datatype=bibtex]{
    %% 容错：把从期刊网页复制的 "et al." 规范成 biblatex 的标准记法 "others"，
    %% 否则它会被当成一位名叫 "al." 的真实作者而显示出来
    \map{\step[fieldsource=author,
                match=\regexp{\bet\s+al\.?},
                replace={others}]}
    \map{\step[fieldsource=sortkey, fieldset=usera, origfieldval]}
  }
}

%% ---------------------------------------------------------------------------
%% 截断词按本人署名位次区分：
%%   本人一作、后面还有作者   ->  C. 7au et al.
%%   本人非一作、后面还有作者 ->  Smith, C. 7au and others
%%   本人已是最后一位作者     ->  Smith and C. 7au（不加后缀）
%% "et al." 是缩写，前面保留逗号；"and others" 是连词短语，前面不加逗号。
%% 不想区分就把下面整块删掉，全文统一用 "et al."
%% ---------------------------------------------------------------------------
\newtoggle{rycfirstauth}
\NewBibliographyString{rycandothers}
\DefineBibliographyStrings{english}{rycandothers={and others}}
\xpatchbibmacro{name:andothers}
  {\bibstring{andothers}}
  {\iftoggle{rycfirstauth}{\bibstring{andothers}}{\bibstring{rycandothers}}}
  {}{}
\xpatchbibmacro{name:andothers}
  {\finalandcomma}
  {\iftoggle{rycfirstauth}{\finalandcomma}{}}
  {}{}

%% 把 maxnames/minnames 设成本人位次 = 打印到本人为止，并记录是否为第一作者
\AtEveryBibitem{%
  \toggletrue{rycfirstauth}%
  \iffieldundef{usera}{}{%
    \defcounter{maxnames}{\thefield{usera}}%
    \defcounter{minnames}{\thefield{usera}}%
    \ifnumgreater{\thefield{usera}}{1}{\togglefalse{rycfirstauth}}{}}}%

\addbibresource{sections/04_publications.bib}


% 字体分支：XeLaTeX/LuaLaTeX 与 pdfLaTeX
\ifxetexorluatex
  \usepackage{fontspec}
  \usepackage[p,osf,swashQ]{cochineal}
  \usepackage[medium,bold]{cabin}
  \usepackage[varqu,varl,scale=0.9]{zi4}
\else
  \usepackage[T1]{fontenc}
  \usepackage[p,osf,swashQ]{cochineal}
  \usepackage{cabin}
  \usepackage[varqu,varl,scale=0.9]{zi4}
\fi

% 可选：页边距 / 配色 / 条目前缀符号，详见 README
% \geometry{left=1cm,right=1cm,top=1.5cm,bottom=1.5cm}
% \definecolor{SwishLineColour}{HTML}{00FFFF}   % 标题下方渐变线
% \definecolor{MarkerColour}{HTML}{0000CC}      % 图标与圆圈编号
% \prefixmarker{$\diamond$}

% 使用说明见 README.md
%
% 载入时机：必须在 \usepackage{settings} 之后（依赖其加载的 biblatex）
% 例外：\PassOptionsToPackage{...}{biblatex} 必须留在 cv-main.tex 里
%       （它得在 settings.sty 载入 biblatex 之前生效）
%

% 去掉 URL 前的 "[Online] Available:" 前缀
\DefineBibliographyStrings{english}{url={\textsc{url}}}

%% 自定义文献排序：年份降序 → 本人作者位次(sortkey)升序 → 作者 → 标题
%% sortkey 由 sections/04_publications.bib 提供，1 = 第一作者
%% 详见 README「参考文献：排序」
\DeclareSortingTemplate{ryc}{
  \sort{\field{presort}}
  \sort[direction=descending]{\field{sortyear}\field{year}\literal{9999}}
  \sort{\field{sortkey}\literal{99}}
  \sort{\field{sortname}\field{author}\field{editor}\field{translator}\field{sorttitle}\field{title}}
  \sort{\field{sorttitle}\field{title}}
}
\ExecuteBibliographyOptions{sorting=ryc}

%% 作者列表截断到本人为止，如 "C. 7au et al."、"Smith, C. 7au and others"
%% sortkey 是 skipout 字段(不写入 .bbl)，故先复制到通用字段 usera 供排版使用
%% 详见 README「作者列表截断」
\DeclareSourcemap{
  \maps[datatype=bibtex]{
    %% 容错：把从期刊网页复制的 "et al." 规范成 biblatex 的标准记法 "others"，
    %% 否则它会被当成一位名叫 "al." 的真实作者而显示出来
    \map{\step[fieldsource=author,
                match=\regexp{\bet\s+al\.?},
                replace={others}]}
    \map{\step[fieldsource=sortkey, fieldset=usera, origfieldval]}
  }
}

%% ---------------------------------------------------------------------------
%% 截断词按本人署名位次区分：
%%   本人一作、后面还有作者   ->  C. 7au et al.
%%   本人非一作、后面还有作者 ->  Smith, C. 7au and others
%%   本人已是最后一位作者     ->  Smith and C. 7au（不加后缀）
%% "et al." 是缩写，前面保留逗号；"and others" 是连词短语，前面不加逗号。
%% 不想区分就把下面整块删掉，全文统一用 "et al."
%% ---------------------------------------------------------------------------
\newtoggle{rycfirstauth}
\NewBibliographyString{rycandothers}
\DefineBibliographyStrings{english}{rycandothers={and others}}
\xpatchbibmacro{name:andothers}
  {\bibstring{andothers}}
  {\iftoggle{rycfirstauth}{\bibstring{andothers}}{\bibstring{rycandothers}}}
  {}{}
\xpatchbibmacro{name:andothers}
  {\finalandcomma}
  {\iftoggle{rycfirstauth}{\finalandcomma}{}}
  {}{}

%% 把 maxnames/minnames 设成本人位次 = 打印到本人为止，并记录是否为第一作者
\AtEveryBibitem{%
  \toggletrue{rycfirstauth}%
  \iffieldundef{usera}{}{%
    \defcounter{maxnames}{\thefield{usera}}%
    \defcounter{minnames}{\thefield{usera}}%
    \ifnumgreater{\thefield{usera}}{1}{\togglefalse{rycfirstauth}}{}}}%

\addbibresource{sections/04_publications.bib}


% 字体分支：XeLaTeX/LuaLaTeX 与 pdfLaTeX
\ifxetexorluatex
  \usepackage{fontspec}
  \usepackage[p,osf,swashQ]{cochineal}
  \usepackage[medium,bold]{cabin}
  \usepackage[varqu,varl,scale=0.9]{zi4}
\else
  \usepackage[T1]{fontenc}
  \usepackage[p,osf,swashQ]{cochineal}
  \usepackage{cabin}
  \usepackage[varqu,varl,scale=0.9]{zi4}
\fi

% 可选：页边距 / 配色 / 条目前缀符号，详见 README
% \geometry{left=1cm,right=1cm,top=1.5cm,bottom=1.5cm}
% \definecolor{SwishLineColour}{HTML}{00FFFF}   % 标题下方渐变线
% \definecolor{MarkerColour}{HTML}{0000CC}      % 图标与圆圈编号
% \prefixmarker{$\diamond$}

% 使用说明见 README.md
%
% 载入时机：必须在 \usepackage{settings} 之后（依赖其加载的 biblatex）
% 例外：\PassOptionsToPackage{...}{biblatex} 必须留在 cv-main.tex 里
%       （它得在 settings.sty 载入 biblatex 之前生效）
%

% 去掉 URL 前的 "[Online] Available:" 前缀
\DefineBibliographyStrings{english}{url={\textsc{url}}}

%% 自定义文献排序：年份降序 → 本人作者位次(sortkey)升序 → 作者 → 标题
%% sortkey 由 sections/04_publications.bib 提供，1 = 第一作者
%% 详见 README「参考文献：排序」
\DeclareSortingTemplate{ryc}{
  \sort{\field{presort}}
  \sort[direction=descending]{\field{sortyear}\field{year}\literal{9999}}
  \sort{\field{sortkey}\literal{99}}
  \sort{\field{sortname}\field{author}\field{editor}\field{translator}\field{sorttitle}\field{title}}
  \sort{\field{sorttitle}\field{title}}
}
\ExecuteBibliographyOptions{sorting=ryc}

%% 作者列表截断到本人为止，如 "C. 7au et al."、"Smith, C. 7au and others"
%% sortkey 是 skipout 字段(不写入 .bbl)，故先复制到通用字段 usera 供排版使用
%% 详见 README「作者列表截断」
\DeclareSourcemap{
  \maps[datatype=bibtex]{
    %% 容错：把从期刊网页复制的 "et al." 规范成 biblatex 的标准记法 "others"，
    %% 否则它会被当成一位名叫 "al." 的真实作者而显示出来
    \map{\step[fieldsource=author,
                match=\regexp{\bet\s+al\.?},
                replace={others}]}
    \map{\step[fieldsource=sortkey, fieldset=usera, origfieldval]}
  }
}

%% ---------------------------------------------------------------------------
%% 截断词按本人署名位次区分：
%%   本人一作、后面还有作者   ->  C. 7au et al.
%%   本人非一作、后面还有作者 ->  Smith, C. 7au and others
%%   本人已是最后一位作者     ->  Smith and C. 7au（不加后缀）
%% "et al." 是缩写，前面保留逗号；"and others" 是连词短语，前面不加逗号。
%% 不想区分就把下面整块删掉，全文统一用 "et al."
%% ---------------------------------------------------------------------------
\newtoggle{rycfirstauth}
\NewBibliographyString{rycandothers}
\DefineBibliographyStrings{english}{rycandothers={and others}}
\xpatchbibmacro{name:andothers}
  {\bibstring{andothers}}
  {\iftoggle{rycfirstauth}{\bibstring{andothers}}{\bibstring{rycandothers}}}
  {}{}
\xpatchbibmacro{name:andothers}
  {\finalandcomma}
  {\iftoggle{rycfirstauth}{\finalandcomma}{}}
  {}{}

%% 把 maxnames/minnames 设成本人位次 = 打印到本人为止，并记录是否为第一作者
\AtEveryBibitem{%
  \toggletrue{rycfirstauth}%
  \iffieldundef{usera}{}{%
    \defcounter{maxnames}{\thefield{usera}}%
    \defcounter{minnames}{\thefield{usera}}%
    \ifnumgreater{\thefield{usera}}{1}{\togglefalse{rycfirstauth}}{}}}%

\addbibresource{sections/04_publications.bib}


% 字体分支：XeLaTeX/LuaLaTeX 与 pdfLaTeX
\ifxetexorluatex
  \usepackage{fontspec}
  \usepackage[p,osf,swashQ]{cochineal}
  \usepackage[medium,bold]{cabin}
  \usepackage[varqu,varl,scale=0.9]{zi4}
\else
  \usepackage[T1]{fontenc}
  \usepackage[p,osf,swashQ]{cochineal}
  \usepackage{cabin}
  \usepackage[varqu,varl,scale=0.9]{zi4}
\fi

% 可选：页边距 / 配色 / 条目前缀符号，详见 README
% \geometry{left=1cm,right=1cm,top=1.5cm,bottom=1.5cm}
% \definecolor{SwishLineColour}{HTML}{00FFFF}   % 标题下方渐变线
% \definecolor{MarkerColour}{HTML}{0000CC}      % 图标与圆圈编号
% \prefixmarker{$\diamond$}

% 使用说明见 README.md
%
% 载入时机：必须在 \usepackage{settings} 之后（依赖其加载的 biblatex）
% 例外：\PassOptionsToPackage{...}{biblatex} 必须留在 cv-main.tex 里
%       （它得在 settings.sty 载入 biblatex 之前生效）
%

% 去掉 URL 前的 "[Online] Available:" 前缀
\DefineBibliographyStrings{english}{url={\textsc{url}}}

%% 自定义文献排序：年份降序 → 本人作者位次(sortkey)升序 → 作者 → 标题
%% sortkey 由 sections/04_publications.bib 提供，1 = 第一作者
%% 详见 README「参考文献：排序」
\DeclareSortingTemplate{ryc}{
  \sort{\field{presort}}
  \sort[direction=descending]{\field{sortyear}\field{year}\literal{9999}}
  \sort{\field{sortkey}\literal{99}}
  \sort{\field{sortname}\field{author}\field{editor}\field{translator}\field{sorttitle}\field{title}}
  \sort{\field{sorttitle}\field{title}}
}
\ExecuteBibliographyOptions{sorting=ryc}

%% 作者列表截断到本人为止，如 "C. 7au et al."、"Smith, C. 7au and others"
%% sortkey 是 skipout 字段(不写入 .bbl)，故先复制到通用字段 usera 供排版使用
%% 详见 README「作者列表截断」
\DeclareSourcemap{
  \maps[datatype=bibtex]{
    %% 容错：把从期刊网页复制的 "et al." 规范成 biblatex 的标准记法 "others"，
    %% 否则它会被当成一位名叫 "al." 的真实作者而显示出来
    \map{\step[fieldsource=author,
                match=\regexp{\bet\s+al\.?},
                replace={others}]}
    \map{\step[fieldsource=sortkey, fieldset=usera, origfieldval]}
  }
}

%% ---------------------------------------------------------------------------
%% 截断词按本人署名位次区分：
%%   本人一作、后面还有作者   ->  C. 7au et al.
%%   本人非一作、后面还有作者 ->  Smith, C. 7au and others
%%   本人已是最后一位作者     ->  Smith and C. 7au（不加后缀）
%% "et al." 是缩写，前面保留逗号；"and others" 是连词短语，前面不加逗号。
%% 不想区分就把下面整块删掉，全文统一用 "et al."
%% ---------------------------------------------------------------------------
\newtoggle{rycfirstauth}
\NewBibliographyString{rycandothers}
\DefineBibliographyStrings{english}{rycandothers={and others}}
\xpatchbibmacro{name:andothers}
  {\bibstring{andothers}}
  {\iftoggle{rycfirstauth}{\bibstring{andothers}}{\bibstring{rycandothers}}}
  {}{}
\xpatchbibmacro{name:andothers}
  {\finalandcomma}
  {\iftoggle{rycfirstauth}{\finalandcomma}{}}
  {}{}

%% 把 maxnames/minnames 设成本人位次 = 打印到本人为止，并记录是否为第一作者
\AtEveryBibitem{%
  \toggletrue{rycfirstauth}%
  \iffieldundef{usera}{}{%
    \defcounter{maxnames}{\thefield{usera}}%
    \defcounter{minnames}{\thefield{usera}}%
    \ifnumgreater{\thefield{usera}}{1}{\togglefalse{rycfirstauth}}{}}}%

\addbibresource{sections/04_publications.bib}


% 字体分支：XeLaTeX/LuaLaTeX 与 pdfLaTeX
\ifxetexorluatex
  \usepackage{fontspec}
  \usepackage[p,osf,swashQ]{cochineal}
  \usepackage[medium,bold]{cabin}
  \usepackage[varqu,varl,scale=0.9]{zi4}
\else
  \usepackage[T1]{fontenc}
  \usepackage[p,osf,swashQ]{cochineal}
  \usepackage{cabin}
  \usepackage[varqu,varl,scale=0.9]{zi4}
\fi

% 可选：页边距 / 配色 / 条目前缀符号，详见 README
% \geometry{left=1cm,right=1cm,top=1.5cm,bottom=1.5cm}
% \definecolor{SwishLineColour}{HTML}{00FFFF}   % 标题下方渐变线
% \definecolor{MarkerColour}{HTML}{0000CC}      % 图标与圆圈编号
% \prefixmarker{$\diamond$}
